\chapter{Complexity — Median}
%%% generated by the blueprint pipeline from TCSlib.Complexity.Randomized.Median
\section{Overview}

Declarations of the Lean module \texttt{TCSlib.Complexity.Randomized.Median}.

\section{Declarations}

\begin{definition}[Characterization of a kth-smallest value]
\lean{Randomized.IsKthSmallest}
\label{Randomized.IsKthSmallest}
\leanok
Let $\alpha$ be a linearly ordered type, let $l$ be a finite list of elements of
$\alpha$, let $k$ be a natural number, and let $y\in\alpha$. The value $y$ is a
kth-smallest value of $l$ if and only if $y$ occurs in $l$, fewer than $k$ entries of
$l$ are strictly less than $y$, and at least $k$ entries are less than or equal to $y$.
\end{definition}

\begin{definition}[Characterization of a median]
\lean{Randomized.IsMedian}
\label{Randomized.IsMedian}
\leanok
Let $\alpha$ be a linearly ordered type, let $l$ be a finite list of elements of
$\alpha$, and let $x\in\alpha$. The element $x$ is a median of $l$ when the length of
$l$ is at most twice the number of entries less than or equal to $x$, and at most twice
the number of entries greater than or equal to $x$.
\end{definition}

\begin{lemma}[Partition of a list at a threshold]
\lean{Randomized.countP_lt_add_countP_le}
\label{Randomized.countP_lt_add_countP_le}
\leanok
\difficulty{4}
Let $\alpha$ be a linearly ordered type, let $l$ be a finite list of elements of
$\alpha$, and let $x\in\alpha$. The number of entries strictly greater than $x$, plus
the number less than or equal to $x$, is the length of $l$.
\end{lemma}

\begin{lemma}[Partition of a list by a threshold]
\lean{Randomized.countP_lt_add_countP_ge}
\label{Randomized.countP_lt_add_countP_ge}
\leanok
\difficulty{4}
Let $\alpha$ be a linearly ordered type, let $l$ be a finite list of elements of
$\alpha$, and let $x\in\alpha$. The number of entries of $l$ strictly less than $x$,
plus the number greater than or equal to $x$, is the length of $l$.
\end{lemma}

\begin{lemma}[Counting entries below or equal to a threshold]
\lean{Randomized.countP_le_eq_countP_lt_add_count}
\label{Randomized.countP_le_eq_countP_lt_add_count}
\leanok
\difficulty{4}
For any list $l$ in a linearly ordered type and any element $y$, the number of entries
at most $y$ is the number strictly less than $y$ plus the number equal to $y$.
\end{lemma}

\begin{theorem}[Characterization of the kth smallest value]
\lean{Randomized.isKthSmallest_iff_of_nodup}
\label{Randomized.isKthSmallest_iff_of_nodup}
\leanok
\uses{Randomized.IsKthSmallest, Randomized.countP_le_eq_countP_lt_add_count}
\difficulty{4}
Let $\alpha$ be a linearly ordered set and let $l$ be a list of distinct elements of
$\alpha$. For a natural number $k$ and $y\in\alpha$, the value $y$ satisfies the
defining rank conditions for being the $k$th smallest element of $l$ if and only if $y$
occurs in $l$ and exactly $k$ entries of $l$ are at most $y$.
\end{theorem}

\begin{theorem}[Uniqueness of an element of fixed rank]
\lean{Randomized.IsKthSmallest.unique}
\label{Randomized.IsKthSmallest.unique}
\leanok
\uses{Randomized.IsKthSmallest}
\difficulty{5}
Let $l$ be a finite list in a linearly ordered set and let $k$ be a natural number. If
$y$ and $z$ both occur in $l$, with fewer than $k$ entries strictly below each and at
least $k$ entries at most each, then $y=z$.
\end{theorem}

\begin{lemma}[Counting below a threshold by a filter]
\lean{Randomized.countP_lt_eq_countP_filter_add}
\label{Randomized.countP_lt_eq_countP_filter_add}
\leanok
\difficulty{4}
Let $\alpha$ be a linearly ordered type, let $l$ be a finite list in $\alpha$, and let
$x,y\in\alpha$ with $x<y$. The number of entries of $l$ strictly less than $y$ equals
the number of entries strictly greater than $x$ and strictly less than $y$, plus the
number of entries less than or equal to $x$.
\end{lemma}

\begin{lemma}[Partition of entries bounded by a threshold]
\lean{Randomized.countP_le_eq_countP_filter_add}
\label{Randomized.countP_le_eq_countP_filter_add}
\leanok
\difficulty{4}
Let $l$ be a list in a linearly ordered type, and let $x<y$. The number of entries of
$l$ at most $y$ equals the number of entries at most $y$ among those strictly greater
than $x$, plus the number of entries at most $x$.
\end{lemma}

\begin{lemma}[Bound on entries below a sorted list element]
\lean{Randomized.countP_lt_getElem_le}
\label{Randomized.countP_lt_getElem_le}
\leanok
\difficulty{5}
Let $\alpha$ be a linearly ordered type, and let $s$ be a finite list in $\alpha$ whose
entries are in nondecreasing order. For any index $j$ with $j<|s|$, the number of
entries of $s$ strictly less than its entry at index $j$ is at most $j$.
\end{lemma}

\begin{lemma}[Bound on larger entries in a sorted list]
\lean{Randomized.countP_getElem_lt_le}
\label{Randomized.countP_getElem_lt_le}
\leanok
\difficulty{5}
Let $s$ be a list in a linearly ordered type, arranged in nondecreasing order, and let
$j$ be an index with $j<|s|$. The number of entries of $s$ strictly greater than the
entry at index $j$ is at most $|s|-1-j$.
\end{lemma}

\begin{lemma}[The kth value of a sorted list]
\lean{Randomized.isKthSmallest_of_sorted}
\label{Randomized.isKthSmallest_of_sorted}
\leanok
\uses{Randomized.IsKthSmallest, Randomized.countP_getElem_lt_le, Randomized.countP_lt_add_countP_le, Randomized.countP_lt_getElem_le}
\difficulty{4}
Let $\alpha$ be a linearly ordered set, and let $s$ and $l$ be lists of elements of
$\alpha$ such that $s$ is a permutation of $l$ and is nondecreasing. Let $k$ be a
natural number with $1\le k$ and $k-1<|s|$. The element at zero-based position $k-1$ in
$s$ satisfies the defining rank conditions for being the $k$th smallest element of $l$:
it occurs in $l$, fewer than $k$ entries are strictly smaller, and at least $k$ entries
are at most it.
\end{lemma}

\begin{theorem}[The kth value of an insertion sort]
\lean{Randomized.isKthSmallest_insertionSort}
\label{Randomized.isKthSmallest_insertionSort}
\leanok
\uses{Randomized.IsKthSmallest, Randomized.isKthSmallest_of_sorted}
\difficulty{1}
Let $\alpha$ be a linearly ordered set, let $l$ be a list of elements of $\alpha$, and
let $k$ be a natural number with $1\le k\le |l|$. The element at zero-based position
$k-1$ in the list obtained by sorting $l$ into nondecreasing order by insertion sort
satisfies the defining rank conditions for being the $k$th smallest element of $l$: it
occurs in $l$, fewer than $k$ entries are strictly smaller, and at least $k$ entries are
at most it.
\end{theorem}

\begin{theorem}[A middle order statistic is a median]
\lean{Randomized.IsKthSmallest.isMedian}
\label{Randomized.IsKthSmallest.isMedian}
\leanok
\uses{Randomized.IsKthSmallest, Randomized.IsMedian, Randomized.countP_lt_add_countP_ge}
\difficulty{3}
Let $l$ be a finite list in a linearly ordered set, and let $y$ be an element whose rank
satisfies $\#\{z\in l:z<y\}<\lfloor(|l|+1)/2\rfloor+1$ and $\lfloor(|l|+1)/2\rfloor+1\le
\#\{z\in l:z\le y\}$. Then at least half the entries of $l$, counting repetitions, are
at most $y$, and at least half are at least $y$; equivalently, $|l|\le 2\#\{z\in l:z\le
y\}$ and $|l|\le 2\#\{z\in l:y\le z\}$.
\end{theorem}

\begin{definition}[Randomized kth-element search]
\lean{Randomized.findKthElement}
\label{Randomized.findKthElement}
\leanok
Let $\alpha$ be a linearly ordered type, let $k$ be a natural number, and let $l$ be a
finite list of elements of $\alpha$. Define a probability distribution on pairs
consisting of an optional element of $\alpha$ and a natural-number cost as follows: for
the empty list it returns no element with cost zero. For a nonempty list, it chooses an
index uniformly at random; if fewer than $k$ entries are strictly below the chosen
pivot, it recursively searches the entries strictly below it, and if fewer than $k$
entries are at most the pivot, it recursively searches the entries strictly above it
with the rank reduced by the number at most the pivot. In either recursive case it adds
the length of the current list to the returned cost; otherwise it returns the pivot with
cost equal to the current list length.
\end{definition}

\begin{theorem}[Correctness of randomized selection]
\lean{Randomized.findKthElement_correct}
\label{Randomized.findKthElement_correct}
\leanok
\uses{Randomized.IsKthSmallest, Randomized.countP_le_eq_countP_filter_add, Randomized.countP_lt_add_countP_le, Randomized.countP_lt_eq_countP_filter_add, Randomized.findKthElement}
\difficulty{6}
Let $\alpha$ be a linearly ordered type, let $l$ be a finite list in $\alpha$, and let
$k$ satisfy $1\le k\le |l|$. Every outcome $(a,c)$ in the support of the randomized
selection procedure on input $(k,l)$ has $a= y$ for some entry $y$ of $l$ such that the
number of entries strictly less than $y$ is less than $k$, and the number of entries
less than or equal to $y$ is at least $k$.
\end{theorem}

\begin{theorem}[Median output of randomized selection]
\lean{Randomized.findKthElement_isMedian}
\label{Randomized.findKthElement_isMedian}
\leanok
\uses{Randomized.IsKthSmallest.isMedian, Randomized.IsMedian, Randomized.findKthElement, Randomized.findKthElement_correct}
\difficulty{4}
Let $\alpha$ be a linearly ordered type and let $l$ be a nonempty finite list in
$\alpha$. Every outcome $(a,c)$ in the support of the randomized selection procedure on
input $(\lfloor(|l|+1)/2\rfloor,l)$ has $a=y$ for some $y\in\alpha$ such that at least
half the entries of $l$, counting multiplicity, are less than or equal to $y$, and at
least half are greater than or equal to $y$.
\end{theorem}

\begin{theorem}[Comparison bound for every outcome]
\lean{Randomized.findKthElement_comparisons_le}
\label{Randomized.findKthElement_comparisons_le}
\leanok
\uses{Randomized.findKthElement}
\difficulty{6}
Let $\alpha$ be a linearly ordered type, let $k$ be a nonnegative integer, and let $l$
be a finite list in $\alpha$. For every outcome $(a,c)$ in the support of the randomized
selection procedure on input $(k,l)$, the number $c$ of comparisons satisfies $2c\le
|l|(|l|+1)$.
\end{theorem}

\begin{definition}[Expected comparison cost]
\lean{Randomized.expectedComparisons}
\label{Randomized.expectedComparisons}
\leanok
\uses{Randomized.findKthElement}
Let $\alpha$ be a linearly ordered type, let $k$ be a natural number, and let $l$ be a
finite list of elements of $\alpha$. The expected comparison cost is the sum, over all
outcomes of the randomized kth-element search, of the probability of that outcome
multiplied by its natural-number cost, regarded as an extended nonnegative real number.
\end{definition}

\begin{definition}[Worst-case expected comparison cost]
\lean{Randomized.maxExpectedComparisons}
\label{Randomized.maxExpectedComparisons}
\leanok
\uses{Randomized.expectedComparisons}
Let $\alpha$ be a linearly ordered type and let $n$ be a natural number. The maximum
expected comparison cost for lists of length $n$ is the supremum of the expected
comparison cost over all natural numbers $k$ and all lists $l$ of elements of $\alpha$
whose length is $n$.
\end{definition}

\begin{definition}[Expected cost of a distribution]
\lean{Randomized.expCost}
\label{Randomized.expCost}
\leanok
Let $\beta$ be a type and let $p$ be a probability distribution on pairs consisting of
an element of $\beta$ and a natural-number cost. Its expected cost is the sum over all
outcomes of the probability of the outcome multiplied by its cost, as an extended
nonnegative real number.
\end{definition}

\begin{lemma}[Expected Cost of a Point Mass]
\lean{Randomized.expCost_pure}
\label{Randomized.expCost_pure}
\leanok
\uses{Randomized.expCost}
\difficulty{1}
For any pair $(b,n)$ consisting of an element $b$ and a natural number $n$, the expected
second coordinate under the point-mass distribution concentrated at $(b,n)$ is $n$.
\end{lemma}

\begin{lemma}[Expected Cost of a Bind]
\lean{Randomized.expCost_bind}
\label{Randomized.expCost_bind}
\leanok
\uses{Randomized.expCost}
\difficulty{3}
Let $p$ be a probability mass function on a type $\gamma$, and for each $a \in \gamma$
let $f(a)$ be a probability mass function on pairs $(b,n)$, where $n$ is a natural
number. The expected second coordinate under the distribution obtained by first sampling
$a$ from $p$ and then sampling from $f(a)$ equals $\sum_a p(a)\,\mathbb{E}_{f(a)}[n]$,
with expectations and sums taken in the nonnegative extended reals.
\end{lemma}

\begin{lemma}[Expected Cost Under an Additive Shift]
\lean{Randomized.expCost_map_add}
\label{Randomized.expCost_map_add}
\leanok
\uses{Randomized.expCost, Randomized.expCost_bind, Randomized.expCost_pure}
\difficulty{3}
Let $p$ be a probability mass function on pairs $(b,m)$, with $m$ a natural number, and
let $n$ be a natural number. If the distribution is changed by replacing each pair
$(b,m)$ with $(b,m+n)$, its expected second coordinate is the expected second coordinate
under $p$ plus $n$.
\end{lemma}

\begin{definition}[Length of the next search interval]
\lean{Randomized.nextLength}
\label{Randomized.nextLength}
\leanok
Let $\alpha$ be a linearly ordered type, let $k$ be a natural number, let $l$ be a
finite list of elements of $\alpha$, and let $x\in\alpha$. Define the next length to be
the number of entries strictly below $x$ if at least $k$ entries are strictly below $x$;
otherwise, define it to be the number of entries strictly above $x$ if fewer than $k$
entries are at most $x$; in all other cases, define it to be zero.
\end{definition}

\begin{theorem}[Expected comparison recurrence]
\lean{Randomized.expectedComparisons_eq}
\label{Randomized.expectedComparisons_eq}
\leanok
\uses{Randomized.expCost, Randomized.expCost_bind, Randomized.expCost_map_add, Randomized.expCost_pure, Randomized.expectedComparisons, Randomized.findKthElement}
\difficulty{4}
Let $\alpha$ be a linearly ordered type, let $k$ be a positive integer, and let $l$ be a
nonempty finite list in $\alpha$. For each position $i$ in $l$, let $a_i$ be the number
of entries of $l$ strictly less than $l_i$, and let $b_i$ be the number of entries less
than or equal to $l_i$. The expected number of comparisons made by the randomized
selection procedure satisfies
\[\mathbb{E}(k,l)=\sum_i \frac{1}{|l|}\left(|l|+\begin{cases}\mathbb{E}(k,\{x\in l:x<l_i\}),&k\le a_i,\\\mathbb{E}(k-b_i,\{x\in l:l_i<x\}),&b_i<k,\\0,&\text{otherwise.}\end{cases}\right),\]
where each filtered list retains the order and multiplicities of its entries.
\end{theorem}

\begin{lemma}[A bound on a sum of maxima]
\lean{Randomized.four_mul_sum_max_le}
\label{Randomized.four_mul_sum_max_le}
\leanok
\difficulty{6}
For every natural number $n$, $4\sum_{j=0}^{n-1}\max\bigl(j,\max(n-1-j,0)\bigr)\le
3n^2$.
\end{lemma}

\begin{lemma}[A bound on list order counts]
\lean{Randomized.four_mul_sum_max_countP_le}
\label{Randomized.four_mul_sum_max_countP_le}
\leanok
\uses{Randomized.countP_getElem_lt_le, Randomized.countP_lt_getElem_le, Randomized.four_mul_sum_max_le}
\difficulty{5}
Let $l$ be a finite list in a linearly ordered type. For each entry $x$ of $l$, count
the entries strictly below $x$ and the entries strictly above $x$. The sum over the
entries of the larger of these two counts, multiplied by $4$, is at most $3|l|^2$, where
$|l|$ is the length of $l$.
\end{lemma}

\begin{theorem}[Linear bound on expected comparisons]
\lean{Randomized.expectedComparisons_le}
\label{Randomized.expectedComparisons_le}
\leanok
\uses{Randomized.expectedComparisons, Randomized.expectedComparisons_eq, Randomized.findKthElement, Randomized.four_mul_sum_max_countP_le, Randomized.nextLength}
\difficulty{6}
Let $\alpha$ be a linearly ordered type, let $k$ be a nonnegative integer, and let $l$
be a finite list in $\alpha$. The expected number of comparisons made by the randomized
selection procedure is at most $10|l|$.
\end{theorem}

\begin{theorem}[Expected comparisons for selection]
\lean{Randomized.maxExpectedComparisons_le}
\label{Randomized.maxExpectedComparisons_le}
\leanok
\uses{Randomized.expectedComparisons_le, Randomized.maxExpectedComparisons}
\difficulty{1}
For every linearly ordered type and every natural number $n$, the supremum of the
expected number of comparisons made by the randomized selection procedure, over all
ranks and all input lists of length $n$, is at most $10n$.
\end{theorem}
